\section{Binary matrices with equal row sums}\label{bin:section}
Curveball trades were introduced for binary matrices with fixed
margins~\cite{SNBFS14}. This section records the statements of
Sections~\ref{edg:section} to~\ref{root:section} in that setting, where the
proofs are the same and the conclusion is different: the comparison leaves
only the functions invariant under row permutations, and on these the bound
can fail when the column sums are unequal. Nothing in this section is used
elsewhere.

Let $m\ge2$ and $k\ge1$, let $1\le r\le k-1$, and let
$c=(c_1,\ldots,c_k)$ be integers with $0<c_j<m$ and $\sum_jc_j=mr$. Let
$\Omega_{r,c}$ be the set of $m\times k$ binary matrices with every row sum
equal to $r$ and column sums $c$, assumed nonempty, with the uniform
measure. For a pair $a=\{x,w\}$ of rows, the $a$-fiber of a matrix $Z$
consists of the matrices in $\Omega_{r,c}$ that agree with $Z$ outside rows
$x$ and $w$. A column is a \emph{singleton} if rows $x$ and $w$ differ in it;
since the two rows have the same sum, there are $2q$ singletons, $q$ with a
one in each row, and the fiber is in bijection with the $q$-subsets of the
singletons. Let $E_a$ average over the $a$-fibers, and put $h_a=I-E_a$,
$H=\sum_ah_a$ over the $\binom m2$ row pairs, and $K=I-H/\binom m2$. Thus
$K$ performs a uniform Curveball trade at a uniformly chosen pair of rows.
The chain is irreducible by Ryser's interchange
theorem~\cite[Theorem~3.1]{Ryser57}, since an interchange is a trade. The
group $S_m$ acts by permuting rows, $L_T$ denotes the sum of $I-T_a$ over the
row transpositions, and a function is \emph{row-symmetric} if it is invariant
under every row permutation.

\begin{proposition}[Row trades]\label{bin:main}
In the setting above:
\begin{enumerate}
\item for every row $x$ and column $j$, the function $\mathbf1[Z_{xj}=1]-c_j/m$
is nonzero and is an eigenfunction of $H$ with eigenvalue $m/2$;
\item $H\succeq\tfrac12L_T$, and $H\succeq\tfrac m2I$ on the orthogonal
complement of the row-symmetric functions;
\item $\operatorname{gap}(H)=\min\{m/2,\mu_{\mathrm{sym}}\}$, where
$\mu_{\mathrm{sym}}$ is the least eigenvalue of $H$ on the centered
row-symmetric functions, and $\mu_{\mathrm{sym}}=\infty$ when there are none.
\end{enumerate}
\end{proposition}

\begin{proof}
(1) As in Lemma~\ref{edg:marginal}, a trade at $\{x,w\}$ gives
$E_a\mathbf1[Z_{xj}=1]=\tfrac12(\mathbf1[Z_{xj}=1]+\mathbf1[Z_{wj}=1])$, and
pairs of rows not containing $x$ fix $\mathbf1[Z_{xj}=1]$. Summing over the
$m-1$ rows $w\ne x$ gives
\[
 H\mathbf1[Z_{xj}=1]=\tfrac{m-1}2\mathbf1[Z_{xj}=1]
 -\tfrac12\bigl(c_j-\mathbf1[Z_{xj}=1]\bigr)
 =\tfrac m2\mathbf1[Z_{xj}=1]-\tfrac{c_j}2 .
\]
By row symmetry the mean of $\mathbf1[Z_{xj}=1]$ is $c_j/m\in(0,1)$, so the
centered function is nonzero and the claim follows from $H\mathbf1=0$.

(2) Exchanging rows $x$ and $w$ preserves $\Omega_{r,c}$ and replaces the set
of singletons with a one in row $x$ by its complement, so it is a trade. The
proof of Proposition~\ref{cmp:main} applies verbatim and gives
$H\succeq\tfrac12L_T$. On the $\lambda$-isotypic component for a partition
$\lambda\ne(m)$ of $m$, with $\lambda_1=m-k'$ and $1\le k'\le m-1$, the
argument of Lemma~\ref{cmp:content} gives
$\operatorname{cont}(\lambda)\le\binom{m-k'}2+\binom{k'}2-k'$, a convex
function of $k'$ whose values at $k'=1$ and $k'=m-1$ are
$\binom{m-1}2-1$ and $\binom{m-1}2-(m-1)$. Hence
$\binom m2-\operatorname{cont}(\lambda)\ge\binom m2-\binom{m-1}2+1=m$, and
$H\succeq\tfrac m2I$ on every component other than $U_{(m)}$, which is the
space of row-symmetric functions.

(3) The space $\mathbf1^\perp$ is the orthogonal sum of the centered
row-symmetric functions and the components $U_\lambda$ with
$\lambda\ne(m)$, each invariant under $H$. On the latter the least eigenvalue
is at least $m/2$ by~(2), and $m/2$ is an eigenvalue on $\mathbf1^\perp$
by~(1). So the least eigenvalue of $H$ on $\mathbf1^\perp$ is the smaller of
$m/2$ and $\mu_{\mathrm{sym}}$.
\end{proof}

For the margins $m=3$, $r=2$, $c=(1,1,2,2)$ the space $\Omega_{r,c}$ has $15$
elements, and $H$ has the eigenvalue $1<3/2$ with a row-symmetric
eigenfunction; this is the computation in the proof of
Proposition~\ref{exp:obstruction}(1), in which the three rows are the
neighborhoods of the vertices of $T$ outside $T$. So equal row sums alone do
not give $\operatorname{gap}(H)=m/2$.

Equal column sums do not give it either. For $m=k=4$ and all row and column
sums equal to $2$, the space has $90$ elements and
$\mu_{\mathrm{sym}}=4/3<2$. For square matrices with all row and column sums
equal to $2$, an exact computation of the same kind as in
Section~\ref{cert:section}, on the classes of matrices under row
permutations, gives $\mu_{\mathrm{sym}}<m/2$ for $m\in\{4,5,6\}$ and
$\mu_{\mathrm{sym}}>m/2$ for $m\in\{7,8\}$; for $m=k=6$ and all sums equal to
$3$ it gives $\mu_{\mathrm{sym}}<3$. The statement corresponding to
Conjecture~\ref{op:conjecture} is therefore false for binary matrices with
equal row sums and equal column sums. In the graph setting rows and columns are permuted by the same relabelings,
so the space $U_{(n)}$ of invariant functions is smaller than the space of
row-symmetric functions of an adjacency matrix.

The expansion of Section~\ref{exp:section} gives a sufficient condition.
Disjoint row pairs resample disjoint rows, so their projections commute and
$D\succeq0$. Let $m\ge3$, and let $\kappa$ be the least eigenvalue, over all triples $T$
of rows, of $H_T$ on the orthogonal complement of its kernel. Then
$H_T^2\succeq\kappa H_T$, and the expansion
$H^2=\sum_TH_T^2-(m-3)H+D$ with $\sum_TH_T=(m-2)H$ gives
\begin{equation}\label{bin:caputo}
 \operatorname{gap}(H)\ge(\kappa-1)(m-2)+1 ,
\end{equation}
which is Caputo's three-particle bound~\cite[Theorem~3.1]{Caputo08} in the
present normalization. Fu, Qin and Wang~\cite[Theorem~5.1]{FQW26} prove
$\kappa\ge1$ for all margins, which gives $\operatorname{gap}(H)\ge1$; the
example with $m=3$ above attains it. If $\kappa\ge3/2$, then
\eqref{bin:caputo} and Proposition~\ref{bin:main}(1) give
$\operatorname{gap}(H)=m/2$. This condition is not necessary: for $m=k=7$
and all sums equal to $2$, fixing the rows $\{3,4\}$, $\{5,6\}$, $\{5,7\}$,
$\{6,7\}$, written as their sets of ones, leaves three rows with the column
sums $(2,2,1,1,0,0,0)$, so $\kappa=1$, while $\mu_{\mathrm{sym}}>7/2$.
